\documentclass[11pt]{article}
\usepackage[margin=0.85in]{geometry}
\usepackage[T1]{fontenc}
\usepackage[utf8]{inputenc}
\usepackage{setspace}
\usepackage{hyperref}
\hypersetup{colorlinks=true,urlcolor=blue,linkcolor=black}
\pagestyle{empty}
\setlength{\parskip}{0.32em}
\setlength{\parindent}{0pt}
\setstretch{0.98}
\begin{document}

\begin{flushleft}
Ke Cheng Hsu \\
Department of Economics \\
The Hong Kong University of Science and Technology \\
Clear Water Bay, Kowloon, Hong Kong \\
\href{mailto:khsu@connect.ust.hk}{khsu@connect.ust.hk} \\
\href{https://cliohsu1997.github.io/}{https://cliohsu1997.github.io/}
\end{flushleft}

\vspace{0.2em}
\today

\vspace{0.35em}
Faculty Search Committee \\
School of Economics \\
Fudan University \\
Shanghai, China

\vspace{0.35em}
Dear Members of the Search Committee,

I am writing to apply for a faculty position at the Assistant Professor level in the School of Economics at Fudan University. I am a Ph.D.\ candidate in Economics at the Hong Kong University of Science and Technology (expected June 2027). My primary research field is Digital Economics, and my research is in digital economics and empirical industrial organization.

My job market paper, ``Designing Return Policy in E-Commerce: A Trade-off between Purchase and Return Efficiency,'' studies how return policy affects both the purchase margin and the return margin. A stricter policy cuts wasteful returns after delivery, but lower return rates make a higher price more profitable for the seller and can price out efficient purchases. Using Taobao apparel data across six product categories, I estimate a structural model of purchase, return, and return-shipping insurance choices. Pairing a discounted price with a moderate restocking fee balances both purchase and return efficiency and performs better than the existing return-shipping insurance policy in most categories.

In my other projects, with Pak Hung Au, I study how heterogeneous search costs sort consumers across niche and standardized products and can raise prices when search costs are lower. With Ce Bian and Liang Zhao, I study pass-through when buyers can return the good and show that a tax can lower the equilibrium price and raise welfare. With Kohei Kawaguchi and Kei Kawai, I study resale price maintenance in Japanese book retail, where fixed retail prices encourage inventory but create unsold stock that is returned and discarded.

At HKUST I teach and assist in undergraduate and postgraduate economics courses in English, including microeconomic theory, information economics, and macroeconomics. My average teaching-assistant evaluation is 4.72/5. As lecturer for Math Camp for the MSc in Economics, my instructor evaluation is 4.53/5. I am prepared to teach microeconomics, industrial organization, information economics, applied microeconometrics, and computational methods for economics.

I have attached my curriculum vitae and job market paper. Three recommendation letters will be sent separately by my referees to this address. Thank you for your consideration. I would be glad to interview at the CICE Job Market in Beijing.

\vspace{0.7em}
Sincerely,

\vspace{0.7em}
Ke Cheng Hsu \\
Ph.D.\ Candidate, Department of Economics, HKUST

\end{document}
